\section{Emparejamiento de flujo condicional (Conditional Flow Matching)}

\subsection{Motivación del establecimiento condicional}

La distribución marginal $p_t(x)$ y el campo vectorial marginal $v_t(x)$ no tienen expresiones analíticas, ya que dependen de la distribución de datos desconocida. Para resolver este problema, el \textbf{emparejamiento de flujo} introduce versiones \textbf{condicionales} de cada magnitud.

La idea central es: dado un punto de datos específico $x_1$, "simplificar" la distribución de datos en una distribución gaussiana concentrada cerca de $x_1$, $\mathcal{N}(x_1, \sigma_{\min}^2 I)$ (donde $\sigma_{\min}$ es extremadamente pequeño). De esta manera, todas las magnitudes condicionales tienen expresiones analíticas.

\begin{knowledgebox}{Establecimiento condicional y proceso de difusión hacia adelante}
Este establecimiento condicional tiene una conexión profunda con el núcleo de transición hacia adelante (forward transition kernel) en los modelos de difusión. En los modelos de difusión, dado un punto de datos $x_0$ (tenga en cuenta la convención de notación de los modelos de difusión), el proceso hacia adelante define la distribución condicional en cada paso de tiempo como $q(x_t | x_0) = \mathcal{N}(\alpha_t x_0, \sigma_t^2 I)$.

En el emparejamiento de flujo, hacemos algo similar: dado un punto de datos $x_1$, definimos la trayectoria de probabilidad condicional $p_t(x | x_1)$. La estructura matemática es completamente idéntica, con la diferencia de que la dirección del tiempo es opuesta.
\end{knowledgebox}

\subsection{Trayectoria de probabilidad condicional}

Dada una muestra de datos $x_1$, la trayectoria de probabilidad condicional se define como:

$$p_t(x | x_1) = \mathcal{N}(x; \mu_t(x_1), \sigma_t^2(x_1) I)$$

donde $\mu_t$ y $\sigma_t$ deben satisfacer las \textbf{condiciones de frontera}:
\begin{itemize}
    \item en $t = 0$: $\mu_0(x_1) = 0$, $\sigma_0(x_1) = 1$, de modo que $p_0(x | x_1) = \mathcal{N}(0, I)$
    \item en $t = 1$: $\mu_1(x_1) \approx x_1$, $\sigma_1(x_1) \approx \sigma_{\min} \to 0$, de modo que $p_1(x | x_1) \approx \delta(x - x_1)$
\end{itemize}

La trayectoria de probabilidad marginal se recupera integrando sobre todos los puntos de datos:

$$p_t(x) = \int p_t(x | x_1) \, p_{\text{data}}(x_1) \, dx_1$$

\subsection{Mapeo de flujo condicional}

\begin{importantbox}{Mapeo de flujo condicional lineal: la belleza de la simplicidad del emparejamiento de flujo}
Una elección de diseño clave del emparejamiento de flujo es definir el mapeo de flujo condicional como una \textbf{función lineal}:
$$\phi_t(x_0 | x_1) = \mu_t(x_1) + \sigma_t(x_1) \cdot x_0$$
Esto significa que, dado un punto de datos $x_1$ y una muestra con ruido $x_0$, la trayectoria desde $x_0$ hasta $x_1$ es una \textbf{línea recta} (cuando $\mu_t$ y $\sigma_t$ son funciones lineales).

Deben satisfacerse las condiciones de frontera:
\begin{itemize}
    \item $\phi_0(x_0 | x_1) = x_0$ (en $t=0$, la salida es la muestra con ruido)
    \item $\phi_1(x_0 | x_1) \approx x_1$ (en $t=1$, la salida se acerca a la muestra de datos)
\end{itemize}
\end{importantbox}

Las ventajas del mapeo de flujo condicional lineal son:
\begin{enumerate}
    \item La trayectoria de probabilidad condicional se convierte naturalmente en una distribución gaussiana
    \item El campo vectorial condicional tiene una expresión analítica
    \item El proceso de muestreo es simple y eficiente
\end{enumerate}

\subsection{Campo vectorial condicional}

Partiendo de la definición del mapeo de flujo condicional, el campo vectorial condicional se puede obtener derivando con respecto al tiempo:

$$u_t(x | x_1) = \frac{d\phi_t(x_0 | x_1)}{dt} = \mu_t'(x_1) + \sigma_t'(x_1) \cdot x_0$$

donde $\mu_t'$ y $\sigma_t'$ son las derivadas de $\mu_t$ y $\sigma_t$ con respecto al tiempo $t$.

Dado que $x_0 = (\phi_t^{-1}(x | x_1)) = \frac{x - \mu_t(x_1)}{\sigma_t(x_1)}$, el campo vectorial condicional se puede escribir como una función de $x$:

$$u_t(x | x_1) = \mu_t'(x_1) + \frac{\sigma_t'(x_1)}{\sigma_t(x_1)} \left( x - \mu_t(x_1) \right)$$

\begin{warningbox}{No confunda las magnitudes condicionales con las marginales}
El campo vectorial condicional $u_t(x | x_1)$ depende del punto de datos específico $x_1$, mientras que finalmente necesitamos el \textbf{campo vectorial marginal} $v_t(x)$, que no depende de $x_1$. El campo vectorial condicional por sí mismo no puede usarse directamente para la generación (ya que durante la generación no conocemos el objetivo $x_1$), pero es una cantidad intermedia clave en el proceso de entrenamiento.
\end{warningbox}

\subsection{De las magnitudes condicionales a las marginales}

El campo vectorial marginal se puede obtener mediante un promedio ponderado del campo vectorial condicional y la trayectoria de probabilidad condicional:

$$v_t(x) = \frac{\int u_t(x | x_1) \, p_t(x | x_1) \, p_{\text{data}}(x_1) \, dx_1}{p_t(x)}$$

\begin{importantbox}{Teorema central del Conditional Flow Matching}
El resultado teórico clave del emparejamiento de flujo es que la función de pérdida utilizada para entrenar con el campo vectorial condicional:
$$\mathcal{L}_{\text{CFM}}(\theta) = \mathbb{E}_{t, \, x_1 \sim p_{\text{data}}, \, x_0 \sim p_0} \left\| v_\theta(t, \phi_t(x_0 | x_1)) - u_t(\phi_t(x_0 | x_1) | x_1) \right\|^2$$
tiene el \textbf{mismo gradiente} (con respecto a los parámetros $\theta$) que la función de pérdida original del emparejamiento de flujo $\mathcal{L}_{\text{FM}}(\theta)$.

Esto significa que podemos entrenar el campo vectorial marginal usando una pérdida condicional, y todas las cantidades en la pérdida condicional tienen expresiones analíticas.
\end{importantbox}

\subsection{Resumen de la sección}

El Conditional Flow Matching convierte magnitudes marginales no analíticas en magnitudes condicionales con forma analítica mediante la introducción de un establecimiento condicional. La garantía teórica clave es que la función de pérdida condicional y la función de pérdida marginal tienen el mismo gradiente, por lo tanto, los efectos del entrenamiento son equivalentes. Esta es la razón fundamental por la cual el emparejamiento de flujo puede evitar la resolución de EDOs.
